% Options for packages loaded elsewhere
\PassOptionsToPackage{unicode}{hyperref}
\PassOptionsToPackage{hyphens}{url}
%
\documentclass[
]{article}
\usepackage{amsmath,amssymb}
\usepackage{lmodern}
\usepackage{iftex}
\ifPDFTeX
  \usepackage[T1]{fontenc}
  \usepackage[utf8]{inputenc}
  \usepackage{textcomp} % provide euro and other symbols
\else % if luatex or xetex
  \usepackage{unicode-math}
  \defaultfontfeatures{Scale=MatchLowercase}
  \defaultfontfeatures[\rmfamily]{Ligatures=TeX,Scale=1}
\fi
% Use upquote if available, for straight quotes in verbatim environments
\IfFileExists{upquote.sty}{\usepackage{upquote}}{}
\IfFileExists{microtype.sty}{% use microtype if available
  \usepackage[]{microtype}
  \UseMicrotypeSet[protrusion]{basicmath} % disable protrusion for tt fonts
}{}
\makeatletter
\@ifundefined{KOMAClassName}{% if non-KOMA class
  \IfFileExists{parskip.sty}{%
    \usepackage{parskip}
  }{% else
    \setlength{\parindent}{0pt}
    \setlength{\parskip}{6pt plus 2pt minus 1pt}}
}{% if KOMA class
  \KOMAoptions{parskip=half}}
\makeatother
\usepackage{xcolor}
\setlength{\emergencystretch}{3em} % prevent overfull lines
\providecommand{\tightlist}{%
  \setlength{\itemsep}{0pt}\setlength{\parskip}{0pt}}
\setcounter{secnumdepth}{-\maxdimen} % remove section numbering
\ifLuaTeX
  \usepackage{selnolig}  % disable illegal ligatures
\fi
\IfFileExists{bookmark.sty}{\usepackage{bookmark}}{\usepackage{hyperref}}
\IfFileExists{xurl.sty}{\usepackage{xurl}}{} % add URL line breaks if available
\urlstyle{same} % disable monospaced font for URLs
\hypersetup{
  hidelinks,
  pdfcreator={LaTeX via pandoc}}

\author{Dietmar Gerald Schrausser}
\date{09.05.2026}

\begin{document}

\hypertarget{folium-irfo.-.1.r}{%
\section{Foliu{[}m{]} Ir/Fo. .1.r}\label{folium-irfo.-.1.r}}

Eine zeitgenössische deutsche Übersetzung wird präsentiert, die auf
einem Vergleich der bestehenden lateinischen, frühneuhochdeutschen
(ENHG) und englischen (Übersetzungs-) Texte basiert, um ein besseres
Verständnis zu ermöglichen. Insbesondere da der ENHG-Text im Vergleich
zum modernen Hochdeutschen schwer verständlich ist, da dieser fundierte
Deutschkenntnisse (als Muttersprache) sowie Kenntnisse verschiedener
Dialekte voraussetzt.

\hypertarget{vorwort-eine-kurze-beschreibung-des-werkes-der-sechs-tage-von-der-erschaffung-der-welt.}{%
\subsection{\texorpdfstring{Vorwort: Eine kurze Beschreibung des
\emph{Werkes der sechs Tage} von der Erschaffung der
Welt.}{Vorwort: Eine kurze Beschreibung des Werkes der sechs Tage von der Erschaffung der Welt.}}\label{vorwort-eine-kurze-beschreibung-des-werkes-der-sechs-tage-von-der-erschaffung-der-welt.}}

Da\footnote{`\textbf{C}Um' und `\textbf{D}Ieweil', vgl. Diodorus
  (\href{https://doi.org/10.3931/e-rara-35702}{1578}, Buch I, pp.5-6).}
unter den gelehrtesten und fähigsten\footnote{`prestantissimos' und
  `fürnamste{[}n{]}'.} Männern, welche die wahre Natur und Geschichte
der (i) Weltentstehung, sowie der (ii) ersten Generation\footnote{`generatione'
  und `Geburt'.} der Menschen beschrieben haben, eine
\emph{zweigeteilte} Meinung herrscht, möchten wir darüber, mit den
\emph{Urzeiten}\footnote{`priscis' und `vordern'.} beginnend und soweit
dies für derart längst vergangene\footnote{`remotis' und
  `enthlegne{[}n{]}'.} Dinge angemessen\footnote{`licebit' und
  `gezime{[}n{]} wil'.} erscheint, in aller Kürze berichten.

Die Ansicht die Welt sei (i) nicht erschaffen worden und sei
unzerstörbar, die Menschheit existiere seit Ewigkeit und habe keinen
Ursprung steht im Gegensatz zur Behauptung, dass die Welt (ii)
geschaffen und vergänglich sei, der Mensch sie seit der ersten
Generation in Besitz genommen habe. Es waren die hochgebildeten
Griechen, welche alle historischen Aufzeichnungen sammelten und der
Theorie folgten, dass vor Beginn sämtlicher Dinge von Himmel und Erde,
solange diese noch vereint waren, lediglich eine einzige Form
existierte. Erst danach nahm die Welt durch die Trennung und Teilung der
Materie die Ordnung und Struktur an, in der wir sie heute sehen.

\begin{enumerate}
\def\labelenumi{(\Roman{enumi})}
\tightlist
\item
  Diese (Griechen) behaupten, dass (i) die Luft und der feurige Teil der
  oberen Schichten ständig in Bewegung waren, leichter wurden und die
  Sonne und die vielen Sterne von diesen getragen werden. Der (ii)
  dunkle und feste Teil jedoch wurde zusammen mit den feuchten Stoffen
  durch sein eigenes Gewicht in die unterste Schicht getragen. Nachdem
  sich diese Stoffe vermischt hatten, entstand aus dem Nebel das Meer,
  die feste, lehmige und weiche Materie wurde zur Erde. Als die Erde
  durch die Hitze der Sonne dichter wurde, bildete sich verrottender
  Schlamm, bedeckt von einer dünnen Haut. Aus solchen Sümpfen und
  Pfützen entstanden vielfältige Lebensformen. Jene, die mehr Wärme
  aufgenommen hatten, wurden zu geflügelten Wesen und stiegen in die
  oberen Schichten auf, die trockeneren und schwereren aber wurden zu
  kriechenden, irdischen Tieren. Solche, die eine wassergebundene Natur
  annahmen, wurden in das Element getragen, das für ihre Art bestimmt
  war.
\end{enumerate}

Als die Erde durch die Hitze der Sonne und die Einwirkung der Luft
austrocknete, entstand eine Mischung vollkommenerer Geschöpfe, männliche
und weibliche. Davon berichtet \textbf{Euripides}, der Tragödiendichter
und \emph{Schüler} des \textbf{Anaxagoras}, des Meisters der
Naturgeschichte. Ebenso sollen die Menschen auf freiem Feld geboren
worden sein, in der Ferne umherziehend, ein wildes und ungeordnetes
Leben führend, welchem Kräuter und Früchte der Bäume Nahrung
boten.\footnote{Anfang des ersten Buches der `Bibliotheca historica' des
  Diodorus Siculus, 1. Jahrhundert v. Chr. (vgl. Diodorus,
  \href{https://doi.org/10.3931/e-rara-68988}{1476}, fol.~15v).}

Nun, da jedoch vieles (Altes wie Neues) über diese Thematik in Latein,
Griechisch und auch Chaldäisch sowie Hebräisch geschrieben wurde, wollen
wir diese alten \emph{fehlerhaften} Darstellungen hinter uns lassen und
uns den \emph{tiefgründigen} mosaischen Schriften über die Schöpfung der
Welt und dem Werk der sechs Tage zuwenden, in welchen die Geheimnisse
der \emph{gesamten} Natur erfasst werden. Denn \textbf{Moses}, der
Prophet und Vater der Geschichtsschreiber Gottes, verstand all dies
durch die \emph{Eingebung} des Heiligen Geistes, des Meisters aller
Wahrheit. Sein menschliches Verständnis sowie seine Erfahrung in
sämtlichen Bereichen des Wissens wurden durch unser Volk, sein Volk und
auch durch die Heiden bezeugt.

König \textbf{Salomon}, ein Erklärer von Natur und Lebewesen, lässt in
seinem Buch der Weisheit erkennen, dass er sein Wissen über diese
wesentlichen Phänomene aus den Gesetzen von Moses bezog. Dieser war (wie
\textbf{Lukas} und \textbf{Philo}, unsere bedeutsamsten Lehrer,
berichten) mit der gesamten ägyptischen Überlieferung bestens vertraut.
Laut \textbf{Hermippus} bezog auch \textbf{Pythagoras} einen Großteil
seiner Philosophie aus dem mosaischen Gesetz.

Der Philosoph \textbf{Numenius} behauptet, dass \textbf{Platon} ein
wahrer attischer Moses gewesen sei, da in den Anfängen seiner Naturlehre
ein Schatz wahrer Weisheit verborgen liegt. Dieser spricht gelehrt und
weise über alle Dinge als von Gott stammend, über ihre Beziehungen
zueinander, ihre Anzahl und die Ordnung sowie ihrer Wandlungen. Es galt
daher bei den alten Hebräern als Gesetz (wie auch \textbf{Hieronymus}
meint), dass man sich nicht mit dieser Schöpfungsgeschichte befassen
sollte, bevor man ein reifes Alter errreicht hat\footnote{`nisi matura
  iam etate atingeret' und `niemant dan der zeitigs alters wer', älter
  als 15 Jahre.}. \emph{Was} nun aber die frommsten Männer,
\textbf{Ambrose} und \textbf{Augustinus}, \textbf{Strabo} auch
\textbf{Beda} oder \textbf{Remigius} und die jüngeren,
\textbf{Aegidius}, \textbf{Albertus}, sowie auch die Griechen
\textbf{Philo}, \textbf{Origenes}, \textbf{Basilius},
\textbf{Theodorus}, \textbf{Apollinarius}, \textbf{Didymus},
\textbf{Gennadius}, \textbf{Chrysostom} usw., über dieses Buch
geschrieben haben, wollen wir unberührt lassen.

Auch werden wir nicht erwähnen, was \textbf{Jonetes}, \textbf{Anchelos}
bzw. \textbf{Simeon der Ältere} in chaldäischer Sprache gesagt haben,
oder was die Hebräer \textbf{Eleazadus}, \textbf{Aba},
\textbf{Johannes}, \textbf{Neonius}, \textbf{Isaak}, \textbf{Josephus},
\textbf{Gersonides}, \textbf{Sadias}, \textbf{Abraham} usw. geschrieben
haben. Vielmehr werden wir kurz die Ordnung der sechs Tage nach Mose
darlegen, in welcher die \emph{göttliche} Erschaffung der Welt
stattfand, was in den heiligen und tiefgründigen religiösen Schriften
erzählt wird.

\begin{enumerate}
\def\labelenumi{(\Roman{enumi})}
\setcounter{enumi}{1}
\tightlist
\item
  Als \emph{Gott} die Welt erschuf, setzte er seinen ersten und größten
  Sohn an die Spitze seines unendlichen Werkes und berief ihn als
  Ratgeber und Meisterhandwerker in die Planung, Verschönerung und
  Gestaltung der Dinge, da dieser mit Weisheit und Verstand im Überfluss
  ausgestattet war.
\end{enumerate}

Jedoch, vor allem mag zu fragen sein, woraus Gott all dieses Große und
Wunderbare, \emph{scheinbar} aus dem Nichts, geschaffen hat. An dieser
Stelle ist es nötig, richterweise die vergänglichen Dinge zu ignorieren
und den Blick auf den \emph{Sitz} zu richten, dorthin, \emph{wo} der
\emph{wahre} Gott wohnt. \emph{Er} verlieh dem Erdreich ewige
Festigkeit, hängte leuchtende Sterne an den Himmel, erschuf die hellsten
Sonnen, umgab die Erde mit dem Meer, ließ die Flüsse fließen, die Felder
sich ausbreiten, die Täler sich senken, die Wälder sich mit Laub
schmücken und die Felsenberge emporragen. All dies hat nicht der
\textbf{Jupiter}\footnote{\begin{enumerate}
  \def\labelenumi{\alph{enumi}.}
  \setcounter{enumi}{18}
  \tightlist
  \item
    Boccaccio
    (\href{https://digi.vatlib.it/view/MSS_Pal.lat.938}{1360}).
  \end{enumerate}}, sondern der Meisterhandwerker der Welt, der
\emph{Quell} des Besten, welcher \emph{Gott} genannt wird und dessen
Anfang unbegreiflich und unergründlich ist, geschaffen.

Die Alten sprachen von drei \emph{Welten}: (i) einer obersten, der Welt
der Engel; (ii) einer himmlischen Welt; und (iii) einer Welt unter dem
Mond, in der wir leben; einer Welt der \emph{Dunkelheit}, die jedoch
regelmäßig vom Licht des Himmels erhellt wird. Darüber hinaus gibt es
eine vierte Welt, in der alle Eigenschaften der anderen Welten vereint
sind: (iv) den Menschen selbst. In der Schule lernten wir, dass der
Mensch eine kleine Welt ist, in der sich durch (i) den aus den Elemente
zusammengesetzten Körper und (ii) einen himmlischen Geist, durch (iii)
die Wachstumsseele der Pflanzen, (iv) die Empfindungsfähigkeit von
vernünftigen Tieren sowie durch (v) die Vernunft und das
engelhafte\footnote{intellektuelle.} Gemüt, das Gleichnis Gottes
erkennen lässt.

Moses berichtet ausführlich von den \emph{drei} Welten und so wie Gott
diese geordnet hat. Auf dem Berg, dort wo er es gelernt hat, wurde ihm
(Moses) auch ausführlich geboten (so wie wir es lesen können), alle
Dinge gemäss dem Gleichnis, welches er dort sah, zu machen. \emph{Was}
nun das Buch Mose über die vollbrachten Werke der sechs Tage lehrt,
werden wir im Überblick darstellen.\footnote{aus dem `Heptaplus'
  (Mirandola \& Mirandola,
  \href{https://doi.org/10.3931/e-rara-61217}{1601}, p.~1 ff.).}

\hypertarget{references}{%
\subsection{References}\label{references}}

Boccaccio, G. (1360). \emph{Genealogiae deorum gentilium libri XV}. Pal.
lat. 938. Vatican: Biblioteca Apostolica Vaticana.
\url{https://digi.vatlib.it/view/MSS_Pal.lat.938}

Diodorus. (1476). \emph{Diodori Sicvli Historiarvm Priscarvm}. Edited by
Poggio Bracciolini, G. F., \& Tacitus, C. Venetiis: per Andreā Iacobi
Katharēsem Andrea Vendramino Duce fortunatissimo.
\url{https://doi.org/10.3931/e-rara-68988}

---------. (1578). \emph{Diodori Siculi Bibliothecae Historicae Libri
XV: \ldots{} Nunc Iterum Latinè \ldots{} Recogniti, \& \ldots{}
Eduntur}. Edited by Grynaeus, J. J., Châteillon, J., Dictys, Dares, \&
Tryphiodorus. Basileae: ex officina Henricpetrina.
\url{https://doi.org/10.3931/e-rara-35702}

Mirandola, G., \& Mirandola, G. F. (1601). \emph{Ioannis Pici,
Mirandulae \ldots{} opera quae extant omnia \ldots{}} Basileae: per
Sebastianum Henricpetri. \url{https://doi.org/10.3931/e-rara-61217}

\end{document}
